\subsection{Task completion can conceal business-state damage}\label{sec:cover-study}
The study compares six cover policies on matched synthetic support tickets, with real LLM decisions and tool execution. It contains 3,600 episodes: two fixed-weight models, six conditions and 300 fixtures per cell, partitioned into 120 training, 60 validation and 120 test fixtures. Each test class has twenty fixtures. The same fixtures and decoy draws are paired across conditions; request phrasings recur across splits.

\paragraph{Execution and observation}
Qwen3:8B and Llama3.1:8B choose checks, business actions, objects, recovery and termination through LangGraph. A controller fetches the ticket and enforces prerequisites without supplying the correct business action. Seven destination-mapped local MCP Streamable HTTP services use actual TLS~1.3/HTTP/2 connections, a common \path{/mcp} path and request/response padding quanta of 4,096/16,384 bytes. The observer receives destinations and bidirectional ciphertext byte counts within supplied episode and call windows. Model traffic and physical timing are outside this observation.

The six conditions are \emph{no cover}, \emph{proxy discard}, \emph{read cover}, \emph{isolated execution}, \emph{live write}, and \emph{isolated execution with visible responses}. A frozen schedule places one or two random decoys among early tool calls. Proxy discard computes a response on a disposable copy without retaining business effects; isolated execution keeps a separate state clone. Live writes modify the primary synthetic fixture. The visible control forwards clearly marked cover responses from the isolated clone to the planner. Other cover responses are hidden. The decoy writes deliberately exercise collateral effects such as changing another order or consuming a shared quota.

\paragraph{Three endpoints on the same runs}
Recognition predicts the requested action even when execution fails. Standardized logistic regression (LR) and a random forest (RF) train within each condition; validation accuracy selects the primary attacker. We also retain both attacks, a frozen baseline-trained attack and destination-position rules. \emph{Goal completion} requires the requested action and object, prerequisites and target postcondition. \emph{State integrity} compares the entire final business state with the single intended-action reference, including other orders, shared resources and ledger entries. A correct final state alone need not satisfy the stricter action/object history required for goal completion.

For each endpoint, five cover-minus-baseline contrasts and one visible-minus-hidden contrast per model form a prespecified twelve-test Holm family. Conditional intervals resample paired test fixtures by action; they do not include learner-selection or deployment uncertainty. All 3,600 formal episodes have terminal results. Supplement S8 gives the frozen design, all attacks, paired comparisons and state replay.

\begin{table*}[t]\centering\small
\caption{Executed cover: all held-out conditions, 120 paired fixtures per model and condition. Recognition columns count correct requested-action predictions; chance is 20/120. The learner is selected using validation LR/RF accuracy. The position rule remains visible when it is stronger. Goal, whole-state integrity and completed tasks with state errors are separate endpoints.}
\label{tab:cover-main}
\setlength{\tabcolsep}{4pt}
\begin{tabularx}{\textwidth}{@{}lYrrrrrr@{}}\toprule
Model & Cover policy & Selected & Last dest. & Goal & State & Goal + error & Call bytes $\times$ \\\midrule
Qwen3:8B & No cover & 118 & 116 & 114 & 116 & 0 & 1.00 \\
Qwen3:8B & Proxy discard & 62 & 77 & 114 & 116 & 0 & 1.34 \\
Qwen3:8B & Read cover & 62 & 77 & 114 & 115 & 0 & 1.34 \\
Qwen3:8B & Isolated, hidden & 61 & 77 & 114 & 115 & 0 & 1.34 \\
Qwen3:8B & Live write & 62 & 75 & 102 & 7 & 95 & 1.34 \\
Qwen3:8B & Isolated, visible & 63 & 77 & 113 & 114 & 1 & 1.35 \\
\midrule
Llama3.1:8B & No cover & 113 & 79 & 82 & 69 & 13 & 1.00 \\
Llama3.1:8B & Proxy discard & 79 & 67 & 80 & 68 & 13 & 1.29 \\
Llama3.1:8B & Read cover & 72 & 67 & 84 & 71 & 14 & 1.28 \\
Llama3.1:8B & Isolated, hidden & 82 & 67 & 83 & 70 & 13 & 1.28 \\
Llama3.1:8B & Live write & 80 & 65 & 76 & 3 & 73 & 1.29 \\
Llama3.1:8B & Isolated, visible & 81 & 68 & 83 & 73 & 11 & 1.29 \\
\bottomrule\end{tabularx}
\end{table*}


\paragraph{Completion masks collateral damage}
Live-write cover completes 102/120 Qwen tasks and 76/120 Llama tasks, but only 7/120 and 3/120 final states match the intended reference (Table~\ref{tab:cover-main}). Thus 95 and 73 completed tasks have incorrect global state. Relative to no cover, integrity falls from 116 to 7 and from 69 to 3: paired differences of $-90.8$ and $-55.0$ percentage points, with conditional 95\% intervals $[-95.0,-85.8]$ and $[-62.5,-47.5]$. Both survive Holm correction. Completion falls by 10.0 and 5.0 points, but only Qwen's contrast is significant (adjusted $p=.022$ versus $1$). Llama already has thirteen completed tasks with state errors without cover, so absolute failures cannot all be attributed to decoys.

\paragraph{Privacy and planner feedback remain separate results}
All ten primary cover-minus-baseline recognition contrasts survive their Holm correction. Live-write scores, 62/120 and 80/120, sit within the ranges observed for the other cover policies. This descriptive overlap does not establish equivalence. The Qwen last-destination rule reaches 75--77/120 under cover and outperforms the selected learner; substantial residual exposure remains. Cover call bytes cost 1.28--1.35 times the paired baseline.

The hidden/visible isolated arms return identical decoy responses in every test pair. Honest action sequences differ in 16 Qwen pairs and 65 Llama pairs, and visible responses increase mean prompt tokens by 18.1\% and 18.9\%. Yet all three primary visible-minus-hidden endpoints have adjusted $p=1$. Hidden arms also exhibit path variation. The control therefore supplies individual execution differences without establishing an average utility or recognition effect.
